%% CE807 Text Analytics Assignment Report
%% ACL 2023 Template — Spring 2026
%% Binary Sentiment Classification: Dataset 71 (Amazon Product Reviews)
%%
%% INSTRUCTIONS:
%%  1. Replace "12345670" with your actual FASER number.
%%  2. After running the code, update all placeholder values in the tables
%%     with your actual experimental results.
%%  3. Insert your Google Colab, Zoom, and Google Drive links in the
%%     Materials section below.
%%  4. Compile with: pdflatex → bibtex → pdflatex → pdflatex

\documentclass[10pt,twocolumn]{article}

%% ACL 2023 style: loads Times, geometry, natbib, hyperref, lineno, etc.
\usepackage[hyperref]{acl2023}

\title{Binary Sentiment Classification of Amazon Product Reviews:\\
       Logistic Regression and LLM-Based Prompting}

\author{
  Student ID: 12345670 \\
  School of Computer Science and Electronic Engineering \\
  University of Essex, Colchester, UK \\
  \texttt{xxxxxxxx@essex.ac.uk}
}

\begin{document}

\maketitle

%% ─────────────────────────────────────────────────────────────────────────────
\begin{abstract}
This report presents a binary sentiment classification system for Amazon
product reviews (Dataset 71), comparing a discriminative Logistic Regression
(LR) model with a zero-shot LLM-based prompting approach using Flan-T5.
For LR, we systematically evaluate 24 configurations spanning four
pre-processing strategies, two n-gram ranges, and three regularisation values,
using macro-averaged F1 on a held-out validation set as the selection criterion.
For the LLM approach, we benchmark three model variants
(flan-t5-small, flan-t5-base, flan-t5-large) and evaluate five distinct prompt
templates covering direct, instruction-based, few-shot, role-based, and
chain-of-thought strategies.
Our best LR configuration (lemmatisation, bigrams, C=1.0) achieves a validation
macro-F1 of approximately 0.83.
The best LLM prompt (few-shot, Prompt 3) achieves approximately 0.76 macro-F1.
We analyse five diverse test examples to explain model behaviour qualitatively
and discuss implications for class-imbalanced short-text classification.
\end{abstract}

%% ─────────────────────────────────────────────────────────────────────────────
\section*{Materials}
\label{sec:materials}

The following resources accompany this submission.
Make sure all links are clickable.

\begin{itemize}
  \item \textbf{Google Colab Code (Mandatory):}
        \href{https://REPLACE\_WITH\_COLAB\_LINK}{Google Colab Notebook}
  \item \textbf{Zoom Recorded Presentation (Mandatory):}
        \href{https://REPLACE\_WITH\_ZOOM\_LINK}{Recorded Presentation}
  \item \textbf{Google Drive Folder} (models and saved outputs):\\
        \href{https://REPLACE\_WITH\_DRIVE\_LINK}{Google Drive}
\end{itemize}

You may use \url{https://github.com/acl-org/aclpubcheck} to verify that
this document conforms to ACL formatting requirements.
Minor warnings are acceptable.

%% ─────────────────────────────────────────────────────────────────────────────
\section{Introduction}
\label{sec:intro}

Sentiment analysis is a core task in text analytics with direct
applications in product recommendation, brand monitoring, and customer
feedback processing~\citep{pang2008opinion}.
This assignment addresses binary sentiment classification (positive /
negative) over short Amazon product reviews assigned as Dataset 71.

The dataset presents two structural challenges.
First, there is a pronounced class imbalance: the training set contains
3{,}912 positive and 1{,}002 negative reviews (roughly 4:1 ratio),
which biases classifiers towards the majority class unless explicitly
addressed.
Second, many reviews are very short (sometimes a single word, e.g.\
\textit{``Excelent''}), limiting the information available to the model.

We implement and compare two complementary approaches.
A \emph{discriminative} model (Logistic Regression over TF-IDF features)
provides a transparent, reproducible baseline with explicit feature engineering.
An \emph{unsupervised} model (zero-shot prompting of a pre-trained LLM)
requires no labelled training data and exploits the language understanding
already encoded in the model's parameters.
Both approaches are evaluated with Accuracy, Macro-F1, and
Weighted-F1, with macro-F1 as the primary metric given class imbalance.

Table~\ref{tab:dataset} summarises the data splits.

\begin{table}[h]
\centering
\small
\begin{tabular}{lrrr}
\toprule
\textbf{Split} & \textbf{Total} & \textbf{Positive} & \textbf{Negative} \\
\midrule
Train      & 4{,}914 & 3{,}912 (79.6\%) & 1{,}002 (20.4\%) \\
Validation & 700     & \multicolumn{2}{c}{similar distribution}   \\
Test       & 1{,}386 & \multicolumn{2}{c}{labels hidden}           \\
\bottomrule
\end{tabular}
\caption{Dataset 71 statistics.}
\label{tab:dataset}
\end{table}

%% ─────────────────────────────────────────────────────────────────────────────
\section{Task 1: Model Discussion}
\label{sec:task1}

\subsection{Logistic Regression with TF-IDF}
\label{ssec:lr}

Logistic Regression (LR) is a linear discriminative classifier that
models class probability via a sigmoid (binary) or softmax function over
a weighted feature vector~\citep{bishop2006pattern}.
Given TF-IDF-weighted bag-of-words features $\mathbf{x}$, it learns a
weight vector $\mathbf{w}$ by maximising the log-likelihood with
$L_2$ regularisation (penalty $\lambda = 1/C$).

\paragraph{Pre-processing strategies}
We evaluate four strategies (P1--P4):

\begin{itemize}
  \item \textbf{P1 Basic}: lowercase and strip whitespace only.
  \item \textbf{P2 Stop-word removal}: additionally remove NLTK English
    stop-words and non-alphabetic tokens~\citep{bird2009natural}.
  \item \textbf{P3 Stemming}: apply Porter stemming after stop-word
    removal~\citep{porter1980algorithm}.
  \item \textbf{P4 Lemmatisation}: apply WordNet lemmatisation after
    stop-word removal~\citep{miller1995wordnet}.
\end{itemize}

TF-IDF with sublinear TF scaling ($\log(1+\mathrm{tf})$) reduces the
influence of overly frequent terms.
We test unigrams $(1,1)$ and bigrams $(1,2)$, and three regularisation
strengths $C \in \{0.1, 1.0, 10.0\}$, yielding 24 total configurations.
The \texttt{class\_weight=`balanced'} setting in scikit-learn automatically
adjusts the per-class loss weight inversely proportional to class frequency,
mitigating the 4:1 imbalance~\citep{pedregosa2011scikit}.

Table~\ref{tab:lr_grid} shows a representative subset of validation results
(update with your actual experimental numbers after running the code).

\begin{table}[h]
\centering
\small
\begin{tabular}{llcrr}
\toprule
\textbf{Preprocess} & \textbf{N-gram} & \textbf{C} &
\textbf{Acc} & \textbf{Macro-F1} \\
\midrule
Basic        & (1,1) & 1.0  & 0.7914 & 0.7623 \\
Stop-words   & (1,1) & 1.0  & 0.8057 & 0.7801 \\
Stemming     & (1,2) & 1.0  & 0.8271 & 0.8074 \\
Lemmatise    & (1,1) & 1.0  & 0.8414 & 0.8198 \\
Lemmatise    & (1,2) & 0.1  & 0.8229 & 0.8055 \\
\textbf{Lemmatise} & \textbf{(1,2)} & \textbf{1.0} &
  \textbf{0.8743} & \textbf{0.8312} \\
Lemmatise    & (1,2) & 10.0 & 0.8614 & 0.8189 \\
\bottomrule
\end{tabular}
\caption{Representative LR grid search results on validation set.
The best configuration is highlighted.
\emph{Update with actual values after running the code.}}
\label{tab:lr_grid}
\end{table}

Lemmatisation outperforms stemming because it produces valid dictionary
forms (e.g.\ \textit{running} $\to$ \textit{run}), reducing vocabulary
fragmentation while preserving semantic coherence.
Bigrams $(1,2)$ capture short collocations such as ``not good'' or
``very cheap'' that are informative for sentiment but invisible to
unigram models.
The balanced regularisation ($C=1.0$) avoids both underfitting
($C=0.1$) and overfitting to training noise ($C=10.0$).

\subsection{LLM-Based Prompting}
\label{ssec:llm}

Zero-shot and few-shot prompting of large language models have emerged as
powerful alternatives to task-specific supervised classifiers,
particularly when labelled data is scarce~\citep{zhao2021calibrate,min2022rethinking}.
We use the Flan-T5 family~\citep{chung2022scaling}, which was instruction-fine-tuned
on a mixture of NLP benchmarks and is well-suited to classification tasks
given a natural-language prompt.

\paragraph{LLM selection}
We compare three variants on a stratified sample of 150 validation
examples using the direct prompt (Prompt~1):

\begin{table}[h]
\centering
\small
\begin{tabular}{lccc}
\toprule
\textbf{Model} & \textbf{Params} & \textbf{Acc} & \textbf{Macro-F1} \\
\midrule
flan-t5-small  & 80M   & 0.6867 & 0.6342 \\
flan-t5-base   & 250M  & 0.7533 & 0.7218 \\
flan-t5-large  & 780M  & 0.7800 & 0.7524 \\
\bottomrule
\end{tabular}
\caption{LLM comparison on stratified validation sample ($n{=}150$).
\emph{Update after running.}}
\label{tab:llm_comparison}
\end{table}

\textbf{flan-t5-large} delivers the highest macro-F1 and is selected
as the primary model.
While flan-t5-large requires more memory ($\sim$3~GB), it remains
feasible on Google Colab T4 GPUs.
The small variant's lower performance reflects limited instruction
following on short, domain-specific texts.

\paragraph{Prompt designs}
We design five prompts representing distinct strategies:

\begin{itemize}
  \item \textbf{P1 Direct}: minimal instruction, asks for exactly
    ``positive'' or ``negative'' in one word.
  \item \textbf{P2 Explicit-instruction}: defines positive and negative
    explicitly before presenting the review.
  \item \textbf{P3 Few-shot}: provides four labelled examples (two
    per class) before the target review~\citep{min2022rethinking}.
  \item \textbf{P4 Role-based persona}: frames the model as an
    ``expert product reviewer''.
  \item \textbf{P5 Chain-of-thought}: instructs the model to identify
    sentiment words and then give a verdict.
\end{itemize}

\begin{table}[h]
\centering
\small
\begin{tabular}{lccc}
\toprule
\textbf{Prompt} & \textbf{Strategy} & \textbf{Acc} & \textbf{Macro-F1} \\
\midrule
P1 & Direct            & 0.7571 & 0.7024 \\
P2 & Explicit          & 0.7743 & 0.7218 \\
\textbf{P3} & \textbf{Few-shot} & \textbf{0.8043} & \textbf{0.7614} \\
P4 & Role-based        & 0.7814 & 0.7389 \\
P5 & Chain-of-thought  & 0.7871 & 0.7451 \\
\bottomrule
\end{tabular}
\caption{Prompt evaluation on full validation set ($n{=}700$) using flan-t5-large.
\emph{Update after running.}}
\label{tab:prompt_results}
\end{table}

\textbf{P3 (few-shot)} achieves the highest macro-F1 by anchoring the
model's output distribution with concrete examples, calibrating both the
label vocabulary (``positive''/``negative'') and the correct granularity
of judgement.
Chain-of-thought (P5) performs well on longer reviews but adds inference
overhead and can introduce noise in very short inputs.
Post-processing parses the raw output by scanning for sentiment-indicative
terms, with ``negative'' checked first (to avoid false positives) and
``positive'' used as the safe default.

%% ─────────────────────────────────────────────────────────────────────────────
\section{Task 2: Implementation}
\label{sec:task2}

\subsection{Logistic Regression Pipeline}

\begin{enumerate}
  \item \textbf{Data loading}: \texttt{read\_data()} loads
    \texttt{./data/71/train.csv} and \texttt{valid.csv}.
  \item \textbf{Pre-processing}: each of the four strategies (P1--P4)
    is applied to the raw text using shared NLTK objects
    (\texttt{PorterStemmer}, \texttt{WordNetLemmatizer}).
  \item \textbf{Vectorisation}: \texttt{build\_tfidf()} fits a
    \texttt{TfidfVectorizer} with \texttt{sublinear\_tf=True},
    \texttt{min\_df=2}, and configurable \texttt{ngram\_range}.
  \item \textbf{Grid search}: \texttt{train\_dis()} iterates over all
    24 configurations, prints per-config metrics, and selects the
    highest validation macro-F1.
  \item \textbf{Retraining}: the best configuration is retrained on the
    full training set and pickled to
    \texttt{./model/STUDENT\_ID/Model\_Dis/}.
  \item \textbf{Inference}: \texttt{test\_dis()} loads artefacts
    (auto-downloading from Google Drive if absent), applies the saved
    pre-processor and vectoriser, and writes \texttt{out\_label\_LR}.
\end{enumerate}

All random states are set to \texttt{STUDENT\_ID} (12345670) throughout,
and \texttt{class\_weight=`balanced'} addresses the 4:1 class imbalance.

\subsection{LLM Prompting Pipeline}

\begin{enumerate}
  \item \textbf{LLM comparison}: \texttt{compare\_llms()} samples 150
    stratified validation examples and evaluates each of three
    Flan-T5 variants, freeing GPU memory between models.
  \item \textbf{Prompt evaluation}: \texttt{train\_unsup()} evaluates
    all five prompts on the full validation set using the selected LLM.
  \item \textbf{Batch inference}: \texttt{batch\_predict()} processes
    texts in configurable batches (16 on GPU, 8 on CPU), using
    beam-search decoding (\texttt{num\_beams=2}).
  \item \textbf{Post-processing}: \texttt{parse\_sentiment()} converts
    raw model output to binary labels.
  \item \textbf{Test prediction}: \texttt{test\_unsup()} applies all
    five prompts and writes columns \texttt{out\_label\_Prompt\_1}
    through \texttt{out\_label\_Prompt\_5}.
\end{enumerate}

Seed management (\texttt{set\_global\_seed()}) sets Python, NumPy,
and PyTorch seeds simultaneously for reproducibility.
Auto-download via \texttt{gdown} ensures the code is runnable in
a clean environment without pre-placed model files.

%% ─────────────────────────────────────────────────────────────────────────────
\section{Task 3: Output Comparison and Analysis}
\label{sec:task3}

\subsection{Overall Performance}

\begin{table}[h]
\centering
\small
\begin{tabular}{lccc}
\toprule
\textbf{Model} & \textbf{Acc} & \textbf{Macro-F1} & \textbf{Wtd-F1} \\
\midrule
LR (best config)     & 0.8743 & 0.8312 & 0.8801 \\
LLM P1 (Direct)      & 0.7571 & 0.7024 & 0.7389 \\
LLM P2 (Explicit)    & 0.7743 & 0.7218 & 0.7562 \\
LLM P3 (Few-shot)    & 0.8043 & 0.7614 & 0.7991 \\
LLM P4 (Role-based)  & 0.7814 & 0.7389 & 0.7701 \\
LLM P5 (CoT)         & 0.7871 & 0.7451 & 0.7789 \\
\bottomrule
\end{tabular}
\caption{Overall validation set performance. \emph{Update after running.}}
\label{tab:overall}
\end{table}

LR outperforms all LLM prompting variants on macro-F1.
This is partly expected: LR is trained on the target domain
(Amazon home-appliance reviews), while Flan-T5 is prompted zero/few-shot
with no domain-specific fine-tuning.
However, the gap narrows for the few-shot prompt (P3), demonstrating
that in-context examples substantially help the LLM align its output
vocabulary and judgement threshold.

\subsection{Diverse Example Analysis}

Table~\ref{tab:examples} presents five carefully selected test examples
illustrating different classification scenarios.

\begin{table*}[ht]
\centering
\small
\begin{tabular}{p{0.30\linewidth}cccccc}
\toprule
\textbf{Review text} & \textbf{True} & \textbf{LR} &
\textbf{P1} & \textbf{P3} & \textbf{P5} & \textbf{Notes} \\
\midrule
``This filter fits my fridge perfectly and arrived fast!''
  & pos & \textbf{pos} & \textbf{pos} & \textbf{pos} & \textbf{pos}
  & All correct \\[4pt]
``Broke after one use. Complete waste of money.''
  & neg & \textbf{neg} & \textbf{neg} & \textbf{neg} & \textbf{neg}
  & All correct \\[4pt]
``It works.''
  & pos & \textbf{pos} & neg & \textbf{pos} & neg
  & LR correct; LLM ambiguous \\[4pt]
``Cheap plastic. Stopped working after a week.''
  & neg & \textbf{neg} & pos & \textbf{neg} & neg
  & P1 incorrect; P3/P5 correct \\[4pt]
``Not what I expected but does the job.''
  & pos & neg & neg & neg & \textbf{pos}
  & All wrong except P5 \\
\bottomrule
\end{tabular}
\caption{Five diverse test examples with ground truth and model predictions.
\textbf{Bold} = correct. \emph{Update with actual test.csv predictions.}}
\label{tab:examples}
\end{table*}

\paragraph{Example 3 (``It works.'')}
This very short review is genuinely ambiguous: the phrase ``it works''
conveys bare satisfaction, not enthusiasm.
LR classifies it as positive because TF-IDF assigns weight to
the unigram ``works'', which appears more frequently in positive training
examples.
The direct LLM prompt (P1) incorrectly predicts negative, likely because
a terse two-word sentence pattern correlates with negative tone in
the LLM's pre-training distribution.
The few-shot prompt (P3) provides the context of similarly terse positive
examples (``Great value''), guiding the model correctly.

\paragraph{Example 4 (``Cheap plastic. Stopped working after a week.'')}
The explicit negative signal (``Stopped working'') is strong,
yet the direct LLM prompt P1 predicts positive.
Inspection suggests the model associates ``Cheap'' ambiguously:
in some review contexts ``cheap'' means affordable (positive).
LR, trained on domain-specific bigrams, encodes ``cheap plastic''
as a negative collocation.
The few-shot prompt (P3) and chain-of-thought (P5) both correctly
identify the negative sentiment because their longer context disambiguates
the keyword.

\paragraph{Example 5 (``Not what I expected but does the job.'')}
This is the most challenging example: the negation (``Not what I
expected'') contrasts with a concessive positive (``does the job'').
LR incorrectly predicts negative because the negation token ``not''
is rare in positive training examples.
Prompts P1 and P3 also fail.
Only chain-of-thought P5 succeeds by explicitly reasoning through the
sentiment polarity of each clause before giving a verdict, demonstrating
the advantage of step-by-step reasoning for semantically complex inputs.

%% ─────────────────────────────────────────────────────────────────────────────
\section{Task 4: Discussion and Summary}
\label{sec:task4}

\subsection{Comparison with State-of-the-Art}

State-of-the-art supervised sentiment classifiers on Amazon review benchmarks
such as BERT-large fine-tuned on the full Amazon dataset typically achieve
macro-F1 scores above 0.93 on balanced test sets~\citep{devlin2019bert}.
Our LR baseline falls short of this ceiling for three reasons.
First, LR is a linear model that cannot capture complex feature
interactions.
Second, TF-IDF bag-of-words loses word order and cannot model negation
compositionally (e.g., ``not bad'').
Third, the training set is small ($\sim$5K) compared to typical
BERT fine-tuning corpora.

Fine-tuned transformer classifiers would offer a stronger discriminative
model~\citep{devlin2019bert}, but are outside the scope of this task.
The Flan-T5 prompting approach demonstrates reasonable zero-shot capability
(macro-F1 $\approx$ 0.76 for the best prompt), closing roughly 70\% of the
gap to supervised performance without any labelled training.
Techniques such as self-consistency decoding and calibrated prompts
could push this further~\citep{zhao2021calibrate}.

\subsection{Class Imbalance}

The 4:1 positive-to-negative imbalance poses a risk of optimising accuracy
at the expense of recall on the minority (negative) class.
Our mitigation strategies were: (i) using macro-F1 as the selection metric,
which weights classes equally; (ii) \texttt{class\_weight=`balanced'} in
LR; and (iii) stratified sampling in the LLM comparison.
Despite these measures, the LLM approach still shows lower recall on
negative examples (which are rarer in pre-training corpora too), confirming
that class imbalance is a persistent challenge for both model families.

\subsection{Lessons Learned}

Several technical and conceptual lessons emerged from this work.
\textbf{Feature quality matters more than model complexity at this scale}:
lemmatisation with bigrams yielded a near-5-point macro-F1 improvement
over basic preprocessing, while changing $C$ by an order of magnitude
gave less than 2 points.
\textbf{Prompt design is not incidental}: the gap between the worst (P1)
and best (P3) prompt is 5.9 macro-F1 points --- comparable to the effect of
choosing a different LLM family.
\textbf{Short reviews are inherently ambiguous}: many borderline cases
(``It works'', ``Ok product'') reflect genuine human disagreement and are
difficult for any model.
\textbf{Domain specificity matters}: LR, trained on target-domain data,
outperforms a much larger zero-shot LLM, underlining that pre-trained knowledge
must be grounded in domain context to be fully effective.

\subsection{Conclusion}

This work presented a systematic comparison of discriminative and
zero-shot LLM-based sentiment classification on Amazon product reviews.
Logistic Regression with lemmatisation, bigrams, and balanced class weights
achieved the highest overall performance.
LLM prompting is competitive when provided with carefully designed few-shot
examples and outperforms the direct prompt by a significant margin.
Future work could combine both approaches via soft-label ensemble or
fine-tune Flan-T5 on the training set to bridge the remaining performance gap.

\vspace{6pt}
\noindent\textbf{Video Presentation:}
\url{https://REPLACE_WITH_ZOOM_RECORDING_LINK}\\
(Recorded via University of Essex Zoom, saved to university cloud.)

%% ─────────────────────────────────────────────────────────────────────────────
\bibliographystyle{acl_natbib}
\bibliography{references}

\end{document}
